\section{Fractional potentials of rank-two forms}
\label{sec:rank-two}

An indefinite rank-two form can have an unbounded density,
but its half-power potential admits a uniform determinant
bound, including for affine terms outside the range of its
quadratic matrix.  Such terms arise under conditioning on
the original coordinates.

\begin{lemma}[A rank-one affine small-ball estimate]
\label{rt:rank-one}
Let \(|u|=1\), \(\lambda\ne0\), and let
\[
 R=\lambda\langle u,X\rangle^2+b^{\mathsf T}X+d.
\]
For every \(w\in\mathbb R\) and \(\rho>0\),
\[
 \mathbb P\{|R-w|<\rho\}
 \leq\min\{1,8\sqrt{\rho/|\lambda|}\}.
\]
\end{lemma}
\begin{proof}
If \(u_i^2\geq1/2\) for some original coordinate, conditioning
on the others leaves a quadratic of leading coefficient
\(\lambda u_i^2\).  Its inverse image of an interval of
length \(2\rho\) has total length at most
\(2\sqrt{2\rho/|\lambda u_i^2|}\leq4\sqrt{\rho/|\lambda|}\).

Otherwise Theorem~\ref{four:diffuse}, with one spectral
direction and the given affine phase, gives
\[
 |\phi_R(t)|\leq\sqrt{e\pi}\,|\lambda t|^{-1/2}.
\]
For \(v>0\), Gaussian Fourier inversion and
\(\boldsymbol1_{\{|z|<\rho\}}\leq e^v e^{-vz^2/\rho^2}\)
imply
\[
 \mathbb P\{|R-w|<\rho\}
 \leq
 \frac{e^v\rho}{2\sqrt{\pi v}}
 \int_{\mathbb R}e^{-\rho^2t^2/(4v)}|\phi_R(t)|\,dt
 \leq
 \frac{e^v\sqrt e\,\Gamma(1/4)}
 {\sqrt2\,v^{1/4}}\sqrt{\rho/|\lambda|}.
\]
Take \(v=1/4\).  The resulting constant
\(e^{3/4}\Gamma(1/4)\) is less than \(8\).
An elementary upper bound sufficient to verify this
numerical inequality is
\[
 \Gamma(1/4)\leq
 4-\frac45+\frac29-\frac2{39}+\frac1{102}+e^{-1}.
\]
Indeed, on \([0,1]\) use the degree-four Taylor upper
bound for \(e^{-x}\), and on \([1,\infty)\) use
\(x^{-3/4}\leq1\).  Together with
\(e^{-1}<3/8\) and \(e^{3/4}<17/8\), the displayed
expression gives a product less than \(8\).
\end{proof}

Define an even function on \(\mathbb R\) by
\[
 \Psi(a)=
 \begin{cases}
 1+\log(1/|a|),&0<|a|\leq1,\\
 |a|^{-1/4},&|a|\geq1,\\
 +\infty,&a=0.
 \end{cases}
\]

\begin{lemma}[A translated quadratic kernel]
\label{rt:kernel}
For \(a\ne0\) and every probability density \(f\leq1\),
\[
 \int_{\mathbb R}f(x)|a-x^2|^{-1/2}\,dx\leq5\Psi(a).
\]
More generally, a probability density bounded by \(K\)
satisfies
\[
 \int_{\mathbb R}f(x)|a-x^2|^{-1/2}\,dx
 \leq5K\Psi(K^2a).
\]
\end{lemma}
\begin{proof}
Let \(m_a(s)\) be the measure of
\(\{x:|a-x^2|^{-1/2}>s\}\).
The integral is at most
\(\int_0^\infty\min\{1,m_a(s)\}\,ds\).
For \(a>0\), put \(\varepsilon=s^{-2}\).
If \(\varepsilon\geq a\), then
\(m_a(s)=2\sqrt{a+\varepsilon}\leq2\sqrt2/s\).
If \(\varepsilon<a\), then
\[
 m_a(s)
 =2\bigl(\sqrt{a+\varepsilon}-\sqrt{a-\varepsilon}\bigr)
 \leq\frac4{\sqrt a\,s^2}.
\]
For \(0<a\leq1\), split the integral at \(1\) and
\(a^{-1/2}\).  It is at most
\(5+\sqrt2\log(1/a)\leq5\Psi(a)\).
For \(a\geq1\), splitting at \(2a^{-1/4}\) gives
at most \(4a^{-1/4}\).

If \(a<0\), then \(m_a(s)\leq2/s\), and it vanishes
when \(s\geq|a|^{-1/2}\).
For \(-1\leq a<0\), split at \(1\) to obtain
\(1+\log(1/|a|)\).  For \(a\leq-1\), the integrand
kernel itself is at most \(|a|^{-1/2}\leq|a|^{-1/4}\).
This proves the first assertion.  For the second,
substitute \(y=Kx\); the rescaled density is bounded
by one and the kernel gains the factor \(K\).
\end{proof}

\begin{lemma}[Averaging the kernel]
\label{rt:kernel-average}
If \(\mathbb P\{|W|<\rho\}\leq L\sqrt\rho\) for all
\(\rho>0\), then
\[
 \mathbb E\Psi(W)\leq4\sqrt L.
\]
\end{lemma}
\begin{proof}
The hypothesis excludes an atom at zero.
For \(0<s<1\), the event \(\Psi(W)>s\) implies
\(|W|<s^{-4}\), so its probability is at most \(Ls^{-2}\).
For \(s\geq1\), it is at most \(Le^{(1-s)/2}\).
Since
\(\sup_{s\geq1}s^2e^{(1-s)/2}=16e^{-3/2}<4\),
both ranges are bounded by \(\min\{1,4L/s^2\}\).
Its integral over \(s>0\) is \(4\sqrt L\).
\end{proof}

\begin{theorem}[Rank-two half-power potential]
\label{rt:half}
Let \(U^{\mathsf T}U=I_2\), let \(\Lambda\) be an
invertible real symmetric \(2\)-by-\(2\) matrix, and
put \(D=|\det\Lambda|\).
For every original-coordinate affine term,
\[
 R=(U^{\mathsf T}X)^{\mathsf T}\Lambda(U^{\mathsf T}X)
       +b^{\mathsf T}X+c_0,
\]
one has
\begin{equation}
 \sup_w\mathbb E|R-w|^{-1/2}
 \leq40\,2^{3/4}D^{-1/4}<68D^{-1/4}.
 \label{eq:rank-two-half}
\end{equation}
\end{theorem}
\begin{proof}
Suppose first that every row leverage of \(U\) is at
most \(1/2\).
Theorem~\ref{four:diffuse} gives
\[
 |\phi_R(t)|\leq\min\{1,K/|t|\},
 \qquad K=e\pi D^{-1/2}.
\]
The Gaussian representation
\[
 |z|^{-1/2}
 =\frac1{\Gamma(1/4)}
   \int_0^\infty v^{-3/4}e^{-vz^2}\,dv
\]
and Gaussian Fourier inversion give
\[
 \mathbb E|R-w|^{-1/2}
 \leq(2\pi)^{-1/2}
       \int_{\mathbb R}|t|^{-1/2}|\phi_R(t)|\,dt.
\]
To justify the order of integration, first write the
Gaussian formula at each \(v>0\), take absolute values,
and then apply Tonelli to the nonnegative integrand.
The integral in \(v\) is
\(\sqrt2\,\Gamma(1/4)|t|^{-1/2}\) before its fixed
prefactor, giving the stated coefficient.
Splitting the last integral at \(|t|=K\) gives
\[
 \mathbb E|R-w|^{-1/2}
 \leq4\sqrt{2e}\,D^{-1/4}.
\]

Otherwise an original coordinate \(i\) has leverage
\(h>1/2\).  Rotate only the two-dimensional spectral
frame so that its row is \((\sqrt h,0)\), and write
\[
 Y_1=\sqrt h\,X_i+T,\qquad Y_2=Z,\qquad
 \Lambda=\begin{pmatrix}\beta&\gamma\\\gamma&\delta\end{pmatrix}.
\]
The variables \(T,Z\) depend on \(X_{-i}\), and \(Z\)
is a unit linear form in those original coordinates.
Put \(k=h^{-1/2}\), \(d=b_i/\sqrt h\), and
\(L_0=b_{-i}^{\mathsf T}X_{-i}-dT+c_0\).
Then
\[
 R=\beta Y_1^2+2\gamma Y_1Z+\delta Z^2+dY_1+L_0.
\]
Conditional on \(X_{-i}\), the density of \(Y_1\) is
bounded by \(k\); its translation is unrestricted.

If \(\beta\ne0\), complete the square:
\[
 R=\beta\left(Y_1+\frac{\gamma Z}{\beta}
                    +\frac d{2\beta}\right)^2+R_0,
 \qquad
 R_0=\frac{\det\Lambda}{\beta}Z^2
       -\frac{\gamma d}{\beta}Z-\frac{d^2}{4\beta}+L_0 .
\]
This is a rank-one quadratic in the remaining original
coordinates plus an arbitrary affine term.
Lemma~\ref{rt:kernel}, conditionally, yields
\[
 \mathbb E[|R-w|^{-1/2}\mid X_{-i}]
 \leq5k|\beta|^{-1/2}\Psi(W),
 \qquad W=k^2(w-R_0)/\beta .
\]
The magnitude of the rank-one coefficient of \(W\) is
\(\lambda_0=k^2D/\beta^2\).
Lemmas~\ref{rt:rank-one} and \ref{rt:kernel-average}
therefore give
\[
 \mathbb E|R-w|^{-1/2}
 \leq20\sqrt8\,k|\beta|^{-1/2}\lambda_0^{-1/4}
 =20\sqrt8\,\sqrt k\,D^{-1/4}
 \leq40\,2^{3/4}D^{-1/4}.
\]
All powers of \(\beta\) cancel.

If \(\beta=0\), then \(D=\gamma^2>0\), and the
conditional polynomial is affine in \(X_i\), of slope
\(\sqrt h(2\gamma Z+d)\).
For any probability density bounded by \(H\),
layer cake gives
\[
 \int f(x)|x-a|^{-1/2}\,dx\leq2\sqrt{2H}.
\]
Apply this first to \(X_i\), then to \(Z\), whose
density is at most \(\sqrt2\) by
Lemma~\ref{four:projection}.  The result is
\[
 \mathbb E|R-w|^{-1/2}
 \leq4\,2^{3/4}h^{-1/4}D^{-1/4}
 \leq8D^{-1/4}.
\]
The exceptional value \(2\gamma Z+d=0\) has probability
zero.  These cases exhaust the possibilities.
\end{proof}

\begin{lemma}[A conditional two-coordinate potential]
\label{rt:conditional}
Suppose \(Z\) has density at most \(H\), and the
conditional density of \(S\) given \(Z\) is at most
\(K\).  For an invertible real symmetric matrix \(L\),
and arbitrary deterministic affine terms,
\[
 \sup_w\mathbb E\left|
 (S,Z)L(S,Z)^{\mathsf T}+\alpha S+\beta Z+\gamma-w
 \right|^{-1/2}
 \leq10\sqrt{15}\sqrt{KH}\,|\det L|^{-1/4}.
\]
\end{lemma}
\begin{proof}
An invertible deterministic quadratic matrix permits
removing its affine term by translating \(S,Z\);
both density caps are preserved.
Write
\(L=\bigl(\begin{smallmatrix}a&b\\b&d\end{smallmatrix}\bigr)\).
For \(a\ne0\), complete the square in \(S\).
Lemma~\ref{rt:kernel} reduces the conditional bound to
\[
 5K|a|^{-1/2}\Psi(A+\epsilon\delta_0Z^2),
 \qquad
 \delta_0=K^2|\det L|/a^2,\quad \epsilon\in\{-1,1\}.
\]
The elementary sublevel estimate
\(|\{z:|z^2+A|<\rho\}|\leq4\sqrt\rho\)
shows that the function
\(v(z)=\Psi(A+\epsilon\delta_0z^2)\) has weak \(L^2\)
norm at most \(\sqrt{15}\delta_0^{-1/4}\).
Indeed, for levels \(0<s<1\),
\(s^2|\{v>s\}|\leq4\delta_0^{-1/2}\).
For \(s\geq1\), the bound is
\(4\delta_0^{-1/2}s^2e^{(1-s)/2}
\leq64e^{-3/2}\delta_0^{-1/2}
<15\delta_0^{-1/2}\).
For any probability density \(g\leq H\), layer cake
gives
\[
 \int gv\leq
 \int_0^\infty\min\{1,H\|v\|_{2,\infty}^2/s^2\}\,ds
 =2\sqrt H\,\|v\|_{2,\infty}.
\]
Substitution cancels \(a\) and proves the asserted
constant.  For \(a=0\), condition first on \(Z\) and
use the affine negative-half-moment estimate from the
last proof; it gives the smaller constant
\(4\sqrt2\sqrt{KH}\,|\det L|^{-1/4}\).
\end{proof}

\begin{corollary}[Affine terms in the range]
\label{rt:range}
In Theorem~\ref{rt:half}, if \(b\in\operatorname{range}(U)\),
then
\[
 \sup_w\mathbb E|R-w|^{-1/2}
 \leq10\sqrt{30}\,D^{-1/4}.
\]
\end{corollary}
\begin{proof}
A deterministic translation of the original coordinates
removes the affine term.  In the diffuse case the
bound \(4\sqrt{2e}D^{-1/4}\) obtained above suffices.
In the heavy-coordinate notation of that proof,
the conditional density of \(Y_1\) given \(Z\) is at
most \(h^{-1/2}\): average the conditional bound given
all \(X_{-i}\).  The marginal density of \(Z\) is at
most \(\sqrt2\).
Lemma~\ref{rt:conditional} gives
\[
 10\sqrt{15}(2/h)^{1/4}D^{-1/4}
 \leq10\sqrt{30}\,D^{-1/4}.
\]
\end{proof}

\begin{theorem}[An exact-rank-three Schur estimate]
\label{rt:rank-three-schur}
Suppose an original-coordinate principal block has
representation
\[
 A_{II}=VMV^{\mathsf T},\qquad
 V^{\mathsf T}V=I_3,\qquad \det M\ne0.
\]
Choose \(i\in I\) with \(a=A_{ii}\ne0\), and write
\(h=|V^{\mathsf T}e_i|^2\).
Then the full ambient quadratic form, with arbitrary
other blocks and affine terms, satisfies
\[
 \|p_Q\|_\infty
 \leq68\bigl(h|a|\,|\det M|\bigr)^{-1/4}.
\]
For a full exact-rank-three matrix with all nonzero
eigenvalues of modulus at least one, this implies
\(\|p_Q\|_\infty\leq68/\sqrt{|a|}\).
\end{theorem}
\begin{proof}
Condition first on the original coordinates outside
\(I\).  They create only an arbitrary affine term in
the independent coordinates of \(I\).
Completing the square in \(X_i\) gives a remainder
with Schur matrix
\[
 S=A_{I\setminus\{i\},I\setminus\{i\}}
 -a^{-1}A_{I\setminus\{i\},i}A_{i,I\setminus\{i\}}.
\]
The square kernel gives
\[
 p_Q(t)\leq |a|^{-1/2}\mathbb E|t-R|^{-1/2}.
\]
Put \(v=V^{\mathsf T}e_i\), choose a \(3\)-by-\(2\)
orthonormal frame \(T\) for \(v^\perp\), and let
\(V_-\) be \(V\) with row \(i\) removed.
Then \(W=V_-T\) satisfies
\[
 W^{\mathsf T}W=T^{\mathsf T}(I-vv^{\mathsf T})T=I_2.
\]
The symmetric matrix
\(H=M-Mvv^{\mathsf T}M/a\) has kernel
\(\operatorname{span}(v)\).  Writing
\(N=T^{\mathsf T}HT\) gives \(S=WNW^{\mathsf T}\).
In a spectral-space basis with \(v=\sqrt h\,e_1\),
write
\[
 M=\begin{pmatrix}\alpha&q^{\mathsf T}\\q&G\end{pmatrix},
 \qquad a=h\alpha.
\]
Then \(N=G-qq^{\mathsf T}/\alpha\), up to an
orthogonal change of its two-dimensional basis, and
\[
 \det N=\det M/\alpha=h\det M/a.
\]
Theorem~\ref{rt:half}, with arbitrary induced affine
terms allowed, proves the first bound.
In the normalized full-rank-three case,
\(|\det M|\geq\|M\|\) and \(|a|\leq h\|M\|\);
hence \(h|a||\det M|\geq |a|^2\), proving the last
assertion.
\end{proof}
